% Shared preamble for build_pdf.py (questions and mark scheme).
% Compiled with Tectonic (XeTeX), so the bank's Unicode text goes in as is.
\documentclass[11pt,a4paper]{article}
\usepackage[a4paper,top=18mm,bottom=20mm,left=24mm,right=18mm]{geometry}
\usepackage{amsmath,amssymb}
\usepackage{fontspec}     % after amsmath; Latin Modern by default
\usepackage{cancel}       % \cancel in mark schemes
\usepackage{pifont}       % \ding for the cross mark
\usepackage{graphicx}
\usepackage{xcolor}
\usepackage{needspace}
\usepackage{newunicodechar}
\usepackage[normalem]{ulem} % \uline: underlined text that can still break

\setlength{\parindent}{0pt}
\setlength{\parskip}{4pt}
\pagestyle{plain}
\tracinglostchars=2        % log every character the font has no glyph for
\emergencystretch=3em      % long maths in mark schemes: stretch spaces rather than overrun
\raggedright               % no wide gaps between words when a line must break early

% Maths wider than the line. A display, or inline maths too wide to fit on a
% line of its own, is scaled down to the line width. Narrower inline maths
% is left alone, so TeX can still break it after = or +.
\newsavebox{\fitbox}
\newcommand{\fitdisplay}[1]{%
  \sbox{\fitbox}{$\displaystyle #1$}%
  \ifdim\wd\fitbox>\linewidth
    \[\resizebox{\linewidth}{!}{\usebox{\fitbox}}\]%
  \else
    \[#1\]%
  \fi}
% Inline maths is scaled a little narrower than the line, leaving room for
% punctuation after it (e.g. "... = 0 *.").
\newcommand{\fitinline}[1]{%
  \sbox{\fitbox}{$#1$}%
  \ifdim\wd\fitbox>\linewidth
    \resizebox{\dimexpr\linewidth-2em\relax}{!}{\usebox{\fitbox}}%
  \else
    $#1$%
  \fi}

% The bank is written for KaTeX, which accepts these Unicode characters in
% text and maths. Latin Modern has no glyph for some of them, so map each one
% to a command that works in both modes.
\newunicodechar{✓}{\ensuremath{\checkmark}}
\newunicodechar{✗}{\ding{55}}
\newunicodechar{≈}{\ensuremath{\approx}}
\newunicodechar{≡}{\ensuremath{\equiv}}
\newunicodechar{≠}{\ensuremath{\neq}}
\newunicodechar{≤}{\ensuremath{\leqslant}}
\newunicodechar{≥}{\ensuremath{\geqslant}}
\newunicodechar{⩽}{\ensuremath{\leqslant}}
\newunicodechar{⩾}{\ensuremath{\geqslant}}
\newunicodechar{ℝ}{\ensuremath{\mathbb{R}}}
\newunicodechar{−}{\ensuremath{-}}
\newunicodechar{×}{\ensuremath{\times}}
\newunicodechar{÷}{\ensuremath{\div}}
\newunicodechar{′}{\ensuremath{'}}
\newunicodechar{•}{\textbullet}

% Question id in small print, so marked work can be matched to bank ids.
\newcommand{\qid}[1]{{\color{gray}\footnotesize\ttfamily #1}}

% Marks in brackets at the right (\qmarks, since \marks is an e-TeX primitive) of the last line, or on a line of their own
% if the last line is full (the TeXbook's right-flush trick).
\newcommand{\qmarks}[1]{{\unskip\nobreak\hfil\penalty50\hskip1.5em\hbox{}\nobreak\hfill
  \textbf{(#1)}\parfillskip=0pt\finalhyphendemerits=0\par}}

% \bankquestion{number}{id}{marks}{working space}{body}
% The body is typeset into a box first to measure it, so the question and its
% working space start on a new page when they would not fit on this one.
% When both together are taller than a page (a big diagram, say), the space
% is cut to what fits beside the question on one page, but never below
% \qminspace; below that it keeps its full size and runs on to the next page.
% \unvbox then puts the body back as ordinary lines, so a question taller
% than a page can still break across pages.
\newlength{\qneed}
\newlength{\qspace}
\newlength{\qminspace}
\setlength{\qminspace}{6cm}
\newcommand{\bankquestion}[5]{%
  \setbox0=\vbox{%
    \noindent\llap{\textbf{\large #1}\hspace{0.8em}}\qid{#2}\par
    #5\qmarks{#3}}%
  \setlength{\qspace}{#4}%
  \ifdim\dimexpr\ht0+\dp0+\qspace\relax>\textheight
    \setlength{\qspace}{\dimexpr\textheight-\ht0-\dp0-3\baselineskip\relax}%
    \ifdim\qspace<\qminspace\setlength{\qspace}{#4}\fi
  \fi
  % If cutting the space by up to a quarter lets the question stay on this
  % page (rather than leave the rest of the page blank), cut it to fit.
  % \pagegoal is \maxdimen on an empty page, which needs no cut.
  \par
  \ifdim\pagegoal<\maxdimen
    \dimen2=\dimexpr\pagegoal-\pagetotal-\ht0-\dp0-2\baselineskip\relax
    \ifdim\dimen2<\qspace\ifdim\dimen2>0.75\qspace\setlength{\qspace}{\dimen2}\fi\fi
  \fi
  \setlength{\qneed}{\dimexpr\ht0+\dp0+\qspace\relax}%
  \ifdim\qneed>\textheight\setlength{\qneed}{\dimexpr\ht0+\dp0\relax}\fi
  \ifdim\qneed<\textheight\Needspace*{\qneed}\fi
  \unvbox0
  % A break is allowed here, so space that must run on to the next page
  % doesn't drag the question's last line with it.
  \penalty0
  \vspace*{\qspace}%
  {\color{lightgray}\hrule}\bigskip}

% \bankfigurepage{image file}{note}: the source page for a diagram question.
\newcommand{\bankfigurepage}[2]{%
  \begin{center}
    \fbox{\includegraphics[width=0.95\linewidth,height=0.55\textheight,keepaspectratio]{#1}}
  \end{center}
  \textit{#2}}

% Mark scheme: \msquestion{number}{id}{marks}, then \mspoint{code}{text}.
\newcommand{\msquestion}[3]{%
  \par\Needspace*{5\baselineskip}\bigskip
  \noindent\textbf{\large #1}\quad\qid{#2}\hfill\textbf{(#3)}\par
  {\color{lightgray}\hrule}\smallskip}
% Both indent through a list rather than \hangindent, so \linewidth (which
% \fitinline scales to) is the real width of the indented lines.
\newenvironment{msindent}[1]{%
  \list{}{\leftmargin=2em\itemindent=#1\labelwidth=0pt\labelsep=0pt
    \topsep=0pt\partopsep=0pt\parsep=\parskip\itemsep=0pt}\item\relax}{\endlist}
\newcommand{\mspoint}[2]{\begin{msindent}{-2em}\textbf{#1}\ #2\end{msindent}}
\newcommand{\msnotes}[1]{\begin{msindent}{0pt}{\small\color{darkgray}#1}\end{msindent}}
\newcommand{\msalt}{\par\smallskip\noindent\textit{Alternative method(s)}\par}
